\section*{Chapitre \ref{ch:b1:complex} --- Nombres complexes}

\begin{solution}{exo:b1:complex:1}
$1 + \iu = \sqrt 2\, \eu^{\iu\pi/4}$~;
$\;\sqrt 3 - \iu = 2\, \eu^{-\iu\pi/6}$~;
$\;-5 = 5\, \eu^{\iu\pi}$~;
\[
\frac{1 + \iu}{\sqrt 3 - \iu}
= \frac{\sqrt 2}{2}\, \eu^{\iu(\pi/4 + \pi/6)}
= \frac{\sqrt 2}{2}\, \eu^{5\iu\pi/12}.
\]
En calculant le même quotient de façon algébrique (multiplication par
le \omterm{def:b1:complex:field}{conjugué})~:
\[
\frac{(1 + \iu)(\sqrt 3 + \iu)}{4}
= \frac{(\sqrt 3 - 1) + \iu(\sqrt 3 + 1)}{4}.
\]
En identifiant avec $\frac{\sqrt 2}{2}(\cos\frac{5\pi}{12} +
\iu\sin\frac{5\pi}{12})$~:
\[
\cos\frac{5\pi}{12} = \frac{\sqrt 6 - \sqrt 2}{4},
\qquad
\sin\frac{5\pi}{12} = \frac{\sqrt 6 + \sqrt 2}{4}.
\]
\end{solution}

\begin{solution}{exo:b1:complex:2}
$(1+\iu)^2 = 2\iu$, donc $(1+\iu)^{20} = (2\iu)^{10} = 2^{10}\,\iu^{10}
= 1024 \times (-1) = -1024$.

Sous forme exponentielle, $(1+\iu)^n = 2^{n/2}\, \eu^{\iu n\pi/4}$, qui
est réel si et seulement si $\sin\frac{n\pi}{4} = 0$, c'est-à-dire $4
\mid n$. Ainsi $(1+\iu)^n \in \R$ exactement pour les multiples de $4$
(la valeur est alors $(-4)^{n/4}$).
\end{solution}

\begin{solution}{exo:b1:complex:3}
$\abs{z - 2} = \abs{z + \iu}$~: équidistance aux points $2$ et $-\iu$
--- c'est la médiatrice du segment joignant $(2, 0)$ et $(0, -1)$.

$\abs{z - 1} = 2$~: le cercle de centre $1$ et de rayon $2$.

$\frac{z-1}{z+1} \in \iu\R$~: d'après la
\cref{prop:b1:complex:geometry}~(3), les points $M(z)$, $A(1)$,
$B(-1)$ vérifient~: les droites $MA$ et $MB$ sont perpendiculaires (ou
bien $z = 1$, où le quotient vaut $0 \in \iu\R$). L'\omterm{def:b1:logic:sets}{ensemble} est le
cercle de diamètre $[-1, 1]$ (le cercle unité), privé du point $-1$ où
le quotient n'est pas défini. \emph{Vérification par le calcul~:} $z =
\eu^{\iu\theta}$ donne $\frac{z-1}{z+1} =
\frac{\eu^{\iu\theta/2}(\eu^{\iu\theta/2} -
\eu^{-\iu\theta/2})}{\eu^{\iu\theta/2}(\eu^{\iu\theta/2} +
\eu^{-\iu\theta/2})} = \iu\tan\frac\theta2 \in \iu\R$.
\end{solution}

\begin{solution}{exo:b1:complex:4}
Linéarisation~:
\[
\sin^4\theta
= \Bigl(\frac{\eu^{\iu\theta} - \eu^{-\iu\theta}}{2\iu}\Bigr)^{\!4}
= \frac{\eu^{4\iu\theta} - 4\eu^{2\iu\theta} + 6 - 4\eu^{-2\iu\theta}
+ \eu^{-4\iu\theta}}{16}
= \frac{\cos 4\theta - 4\cos 2\theta + 3}{8}.
\]
Développement~: par de Moivre, $\cos 4\theta = \Re\bigl((c +
\iu s)^4\bigr) = c^4 - 6c^2 s^2 + s^4$ avec $c = \cos\theta$, $s =
\sin\theta$~; en substituant $s^2 = 1 - c^2$~:
\[
\cos 4\theta = c^4 - 6c^2(1 - c^2) + (1 - c^2)^2
= 8c^4 - 8c^2 + 1 .
\]
\end{solution}

\begin{solution}{exo:b1:complex:5}
$8\iu = 8\,\eu^{\iu\pi/2}$, donc les racines cubiques sont
$2\,\eu^{\iu(\pi/6 + 2k\pi/3)}$, $k = 0, 1, 2$~:
\[
z_0 = 2\eu^{\iu\pi/6} = \sqrt 3 + \iu,\quad
z_1 = 2\eu^{5\iu\pi/6} = -\sqrt 3 + \iu,\quad
z_2 = 2\eu^{3\iu\pi/2} = -2\iu :
\]
un triangle équilatéral sur le cercle de rayon $2$.

$-4 = 4\eu^{\iu\pi}$, donc $z^4 = -4$ a pour solutions $\sqrt 2\,
\eu^{\iu(\pi/4 + k\pi/2)}$~: $\;1 + \iu$, $-1 + \iu$, $-1 - \iu$,
$1 - \iu$. En appariant les racines \omterm{def:b1:complex:field}{conjuguées}~:
\[
X^4 + 4 = \bigl(X^2 - 2X + 2\bigr)\bigl(X^2 + 2X + 2\bigr),
\]
puisque $(X - (1+\iu))(X - (1-\iu)) = X^2 - 2X + 2$, et de même pour
l'autre paire. (Développer pour vérifier.)
\end{solution}

\begin{solution}{exo:b1:complex:6}
$C_n + \iu S_n = \sum_{k=0}^{n} \eu^{\iu k\theta} =
\frac{\eu^{\iu(n+1)\theta} - 1}{\eu^{\iu\theta} - 1}$ (somme
géométrique, de raison $\eu^{\iu\theta} \neq 1$). Factorisation par
l'angle moitié (\cref{met:b1:complex:trig}~(4)) au numérateur et au
dénominateur~:
\[
\frac{2\iu\sin\frac{(n+1)\theta}{2}\, \eu^{\iu(n+1)\theta/2}}
{2\iu\sin\frac{\theta}{2}\, \eu^{\iu\theta/2}}
= \frac{\sin\frac{(n+1)\theta}{2}}{\sin\frac{\theta}{2}}\,
\eu^{\iu n\theta/2}.
\]
En prenant les parties réelle et imaginaire~:
\[
C_n = \frac{\sin\frac{(n+1)\theta}{2}}{\sin\frac{\theta}{2}}
\cos\frac{n\theta}{2},
\qquad
S_n = \frac{\sin\frac{(n+1)\theta}{2}}{\sin\frac{\theta}{2}}
\sin\frac{n\theta}{2}.
\]
\end{solution}

\begin{solution}{exo:b1:complex:7}
Discriminant~: $\Delta = (3 + 4\iu)^2 - 4(-1 + 5\iu) = 9 + 24\iu - 16 +
4 - 20\iu = -3 + 4\iu$. Racines carrées de $-3 + 4\iu$~: on résout $x^2
- y^2 = -3$, $2xy = 4$, $x^2 + y^2 = 5$~; alors $x^2 = 1$, $y = 2/x$,
de même signe~: $\delta = \pm(1 + 2\iu)$. D'où
\[
z = \frac{(3 + 4\iu) \pm (1 + 2\iu)}{2}
\in \{\, 2 + 3\iu,\; 1 + \iu \,\}.
\]
\emph{Vérification~:} la somme vaut $3 + 4\iu$ et le produit
$(2+3\iu)(1+\iu) = -1 + 5\iu$, comme l'exigent les coefficients.
\end{solution}

\begin{solution}{exo:b1:complex:8}
\begin{enumerate}
  \item \cref{prop:b1:complex:sumroots} avec $n = 5$.
  \item $u + v = \omega + \omega^2 + \omega^3 + \omega^4 = -1$ d'après
        (1). Pour le produit, on développe et on réduit les exposants
        modulo $5$~:
        \[
        uv = (\omega + \omega^4)(\omega^2 + \omega^3)
        = \omega^3 + \omega^4 + \omega^6 + \omega^7
        = \omega^3 + \omega^4 + \omega + \omega^2 = -1 .
        \]
        Ainsi $u$ et $v$ ont pour somme $-1$ et pour produit $-1$~: ce
        sont les deux racines de $X^2 + X - 1$.
  \item $u = \omega + \conj\omega = 2\cos\frac{2\pi}{5} > 0$ (l'angle
        est aigu), et la racine positive de $X^2 + X - 1$ est
        $\frac{-1 + \sqrt 5}{2}$. D'où $\cos\frac{2\pi}{5} =
        \frac{\sqrt 5 - 1}{4}$.
\end{enumerate}
\end{solution}

\begin{solution}{exo:b1:complex:9}
On développe les deux carrés comme dans la démonstration de la
\cref{prop:b1:complex:rules}~(4)~:
\[
\abs{z + w}^2 = \abs z^2 + \abs w^2 + 2\Re(z\conj w),
\qquad
\abs{z - w}^2 = \abs z^2 + \abs w^2 - 2\Re(z\conj w),
\]
et on ajoute. Géométriquement, $\abs{z+w}$ et $\abs{z-w}$ sont les
longueurs des deux diagonales du parallélogramme de sommets $0, z,
z+w, w$, tandis que $\abs z$ et $\abs w$ sont les longueurs des côtés~:
la somme des carrés des diagonales est égale à la somme des carrés des
quatre côtés.
\end{solution}

\begin{solution}{exo:b1:complex:10}
Posons $j = \eu^{2\iu\pi/3}$, de sorte que $j^2 + j + 1 = 0$ et
$\eu^{\iu\pi/3} = -j^2$, $\eu^{-\iu\pi/3} = -j$. Le triangle est
équilatéral direct si et seulement si $c - a = -j^2 (b - a)$, indirect
si et seulement si $c - a = -j(b - a)$
(\cref{ex:b1:complex:equilateral}).

Considérons $P = a + jb + j^2 c$ et $Q = a + j^2 b + jc$. En utilisant
$1 + j^2 = -j$ et $j^3 = 1$~:
\[
c - a + j^2(b - a) = c + j^2 b - (1 + j^2)\,a = c + j^2 b + ja
= j\,(a + jb + j^2 c) = jP,
\]
si bien que le cas direct s'écrit $jP = 0 \iff P = 0$~; le même calcul
avec $j$ à la place de $j^2$ donne $c - a + j(b - a) = j^2 Q$, si bien
que le cas indirect s'écrit $Q = 0$. Donc~: équilatéral (quelle que
soit l'orientation) $\iff PQ = 0$. En développant, avec $j + j^2 =
-1$~:
\[
PQ = a^2 + b^2 + c^2 + (j + j^2)(ab + bc + ca)
= a^2 + b^2 + c^2 - (ab + bc + ca).
\]
Le triangle est donc équilatéral si et seulement si $a^2 + b^2 + c^2 =
ab + bc + ca$.
\end{solution}

\begin{solution}{exo:b1:complex:11}
Comme les $\omega^k$, $k = 0, \dots, n-1$, sont exactement les $n$
racines de $X^n - 1$ (\cref{thm:b1:complex:roots}) et que le polynôme
est unitaire~:
\[
X^n - 1 = \prod_{k=0}^{n-1} (X - \omega^k)
= (X - 1) \prod_{k=1}^{n-1} (X - \omega^k).
\]
En divisant par $X - 1$~: $\;1 + X + \dots + X^{n-1} =
\prod_{k=1}^{n-1} (X - \omega^k)$. En évaluant en $X = 1$, on obtient
$P = n$.

Or $1 - \omega^k = -\eu^{\iu k\pi/n}\bigl(\eu^{\iu k\pi/n} -
\eu^{-\iu k\pi/n}\bigr) = -2\iu\,\eu^{\iu k\pi/n} \sin\frac{k\pi}{n}$,
donc, en prenant les \omterm{def:b1:complex:field}{modules} dans $P = n$ (chaque $\sin\frac{k\pi}n >
0$ pour $1 \leq k \leq n-1$)~:
\[
n = \abs P = \prod_{k=1}^{n-1} 2\sin\frac{k\pi}{n}
= 2^{n-1} \prod_{k=1}^{n-1} \sin\frac{k\pi}{n},
\qquad\text{d'où}\qquad
\prod_{k=1}^{n-1} \sin\frac{k\pi}{n} = \frac{n}{2^{n-1}} .
\]
\end{solution}

\begin{solution}{exo:b1:complex:12}
$z = 1$ n'est pas solution ($2^n \neq 0$), de sorte que l'équation
équivaut à $\bigl(\frac{z+1}{z-1}\bigr)^n = 1$, c'est-à-dire
$\frac{z+1}{z-1} = \omega^k$ avec $\omega = \eu^{2\iu\pi/n}$ et
$k \in \{0, \dots, n-1\}$. La valeur $k = 0$ est exclue ($z + 1 =
z - 1$ est impossible). Pour $1 \leq k \leq n - 1$, la résolution de
$z + 1 = \omega^k(z - 1)$ donne $z(1 - \omega^k) = -1 - \omega^k$,
donc, avec $\varphi = \frac{2k\pi}{n}$ et les factorisations par
l'angle moitié $1 + \eu^{\iu\varphi} = 2\cos\frac\varphi2\,
\eu^{\iu\varphi/2}$, $\eu^{\iu\varphi} - 1 = 2\iu\sin\frac\varphi2
\,\eu^{\iu\varphi/2}$~:
\[
z_k = \frac{1 + \omega^k}{\omega^k - 1}
= \frac{2\cos\frac{k\pi}{n}}{2\iu\,\sin\frac{k\pi}{n}}
= -\iu\,\frac{\cos\frac{k\pi}{n}}{\sin\frac{k\pi}{n}} ,
\]
imaginaire pur comme annoncé. L'\omterm{def:b1:logic:map}{application} $t \mapsto \cos t/\sin t$
est \omterm{def:b1:logic:inj}{injective} sur $\intoo0\pi$ (elle y est strictement décroissante),
donc les $n - 1$ valeurs $z_k$ sont deux à deux distinctes~:
l'équation, de degré $n - 1$ une fois développée (les termes en $z^n$
se simplifient), a exactement ces $n - 1$ solutions.
\end{solution}

\begin{solution}{pb:b1:complex:1}
\textbf{1.} $j = \cos\frac{2\pi}3 + \iu\sin\frac{2\pi}3 = -\frac12
+ \iu\frac{\sqrt3}2$~; $j^2 = \eu^{4\iu\pi/3} = -\frac12 -
\iu\frac{\sqrt3}2 = \conj j$~; $j^3 = 1$~; $1 + j + j^2 = 0$ (somme
des racines cubiques de l'unité, \cref{prop:b1:complex:sumroots})~;
$j^{-1} = j^2$ (car $j \cdot j^2 = 1$). Sur le cercle unité, $1$,
$j$, $j^2$ sont les sommets d'un triangle équilatéral direct.

\textbf{2.} La rotation de centre $a$ et d'angle $\theta$ fixe $a$
et fait tourner de $\theta$ tout vecteur issu de $a$~: $r(z) - a =
\eu^{\iu\theta}(z - a)$, c'est-à-dire $r(z) = a + \eu^{\iu\theta}(z
- a)$. En composant $r(z) = a + \eu^{\iu\theta}(z-a)$ et $r'(z) = a'
+ \eu^{\iu\theta'}(z-a')$~:
\[
r' \circ r\,(z) = \eu^{\iu(\theta + \theta')} z +
\text{constante},
\]
une \omterm{def:b1:logic:map}{application} de la forme $Az + B$ avec $A =
\eu^{\iu(\theta+\theta')}$ de \omterm{def:b1:complex:field}{module} $1$. D'après la
\cref{prop:b1:complex:similitude}, c'est une rotation d'angle $\arg
A = \theta + \theta'$ si $A \neq 1$, et une translation si $A = 1$,
c'est-à-dire si $\theta + \theta' \in 2\pi\Z$.

\textbf{3.} $(b, c, p)$ est équilatéral si et seulement si $\abs{p -
b} = \abs{c - b} = \abs{p - c}$. En posant $q = \frac{p - b}{c -
b}$, la première égalité dit $\abs q = 1$ et la seconde $\abs{q - 1}
= 1$~; \omterm{def:b1:logic:sets}{ensemble}, $q = \eu^{\pm\iu\pi/3}$ (les deux points
d'intersection des cercles $\abs q = 1$ et $\abs{q - 1} = 1$ sont
$\frac12 \pm \iu\frac{\sqrt3}2$). D'où $p = b + \eu^{\pm\iu\pi/3}(c
- b)$. Pour $a = 0$, $b = 1$, $c = \iu$ (un triangle direct)~: le
choix $\eu^{-\iu\pi/3}$ donne
\[
p = 1 + \Bigl(\tfrac12 - \iu\tfrac{\sqrt3}2\Bigr)(\iu - 1)
= \tfrac{1 + \sqrt3}2\,(1 + \iu) \approx 1{,}37 + 1{,}37\,\iu ,
\]
qui se situe de l'autre côté de la droite $BC$ ($x + y = 1$) par
rapport à $a = 0$~: vers l'extérieur. Le choix $\eu^{+\iu\pi/3}$
donne $p \approx -0{,}37 - 0{,}37\,\iu$, du même côté que $a$~: vers
l'intérieur.

\textbf{4.} Multiplions $P = a + jb + j^2c$ par $j^2$~: $j^2 P = j^2
a + j^3 b + j^4 c = b + jc + j^2 a$ (en utilisant $j^3 = 1$, $j^4 =
j$)~; et par $j$~: $jP = ja + j^2 b + c$. Ce sont les deux
identités. Si $P = 0$, alors $j^2 P = jP = 0$~: le critère vaut
aussi pour $(b, c, a)$ et $(c, a, b)$ --- invariance cyclique (comme
il se doit~: un triangle équilatéral se moque de savoir quel sommet
est cité en premier).

\textbf{5.} Si $a = b$~: $0 = a(1 + j) + j^2 c = -j^2 a + j^2 c$ (en
utilisant $1 + j = -j^2$), donc $c = a$. Si $b = c$~: $0 = a + b(j +
j^2) = a - b$, donc $a = b$. Si $a = c$~: $0 = a(1 + j^2) + jb = -ja
+ jb$, donc $a = b$. Dans chaque cas, les trois points coïncident.

\textbf{6.} $(a'z + b') \circ (az + b) = a'a\,z + (a'b + b')$ avec
$a'a \neq 0$~: même forme. La réciproque de $z \mapsto az + b$ est
$z \mapsto \frac1a z - \frac ba$, encore de la même forme. Avec
l'\omterm{def:b1:logic:map}{application} identité comme élément neutre, les similitudes
directes forment un groupe pour la composition.

\textbf{7.} $f(z) = az + b$ vérifie $f(z_1) = w_1$, $f(z_2) = w_2$
si et seulement si $a z_1 + b = w_1$ et $a z_2 + b = w_2$~; par
soustraction, $a(z_2 - z_1) = w_2 - w_1$, donc
\[
a = \frac{w_2 - w_1}{z_2 - z_1} \;(\neq 0), \qquad
b = w_1 - a z_1
\]
sont imposés, et réciproquement ce choix convient~: existence et
unicité.

\textbf{8.} $\abs{f(z) - f(w)} = \abs{a(z - w)} = \abs a\,\abs{z -
w}$~: toutes les distances sont multipliées par $\abs a$. Pour la
forme~:
\[
\frac{f(c') - f(a')}{f(b') - f(a')} = \frac{a(c' -
a')}{a(b' - a')} = \frac{c' - a'}{b' - a'} .
\]
Si deux triangles $(a', b', c')$ et $(a'', b'', c'')$ ont la même
forme, soit $f$ l'unique similitude directe telle que $f(a') = a''$
et $f(b') = b''$ (question 7)~; alors la forme de $(a'', b'',
f(c'))$ est égale à celle de $(a', b', c')$, donc à celle de $(a'',
b'', c'')$, et la forme détermine le troisième sommet à partir des
deux premiers~: $f(c') = c''$. Réciproquement, l'égalité des formes
résulte de l'invariance affichée ci-dessus.

\textbf{9.} $b = f(0) = 1$~; $a + b = f(1) = \iu$ donne $a = \iu -
1$. Rapport $\abs a = \sqrt2$, angle $\arg(\iu - 1) =
\frac{3\pi}4$. Point fixe~:
\[
\zeta = \frac{b}{1 - a} = \frac1{2 - \iu} = \frac{2 + \iu}5 .
\]
Ainsi $f$ est la similitude directe de centre $\frac{2+\iu}5$, de
rapport $\sqrt2$ et d'angle $\frac{3\pi}4$.

\textbf{10.} Avec $\alpha + \beta = 1$ et $f(z) = az + b$~:
\[
f(\alpha a' + \beta b') = a\alpha a' + a\beta b' + b
= \alpha(a a' + b) + \beta(a b' + b)
= \alpha f(a') + \beta f(b') ,
\]
la clé étant $b = (\alpha + \beta)b$. Les milieux ($\alpha = \beta =
\frac12$), les isobarycentres (en itérant) et les centres de
Napoléon $\alpha b + \beta c$ ci-dessous sont donc
\emph{équivariants}~: transformer le triangle les transforme en
conséquence. Démontrer un énoncé invariant par similitude pour un
seul triangle bien placé le démontre donc pour tous --- c'est la
licence utilisée à la question 21.

\textbf{11.} Le sommet extérieur sur $[b, c]$ est $p_a = b +
\eu^{-\iu\pi/3}(c - b)$ (question 3), donc le centre vaut
\[
n_a = \frac{b + c + p_a}3
= \frac{2b + c + \frac{1 - \iu\sqrt3}2\,(c - b)}3
= \frac{(3 + \iu\sqrt3)\,b + (3 - \iu\sqrt3)\,c}6
= \alpha b + \beta c ,
\]
avec $\alpha = \frac{3 + \iu\sqrt3}6$ et $\beta = \frac{3 -
\iu\sqrt3}6 = \conj\alpha$. Le même calcul sur les côtés $[c, a]$ et
$[a, b]$ donne $n_b = \alpha c + \beta a$ et $n_c = \alpha a + \beta
b$ (décalage cyclique des rôles).

\textbf{12.} $\alpha + \beta = \frac{3 + \iu\sqrt3 + 3 -
\iu\sqrt3}6 = 1$. D'où
\[
\frac{n_a + n_b + n_c}3
= \frac{(\alpha + \beta)(a + b + c)}3 = \frac{a + b + c}3 :
\]
les deux triangles ont le même isobarycentre.

\textbf{13.} On regroupe selon $\alpha$ et $\beta$ et on applique
les identités de la question 4 à $P = a + jb + j^2c$~:
\[
n_a + j n_b + j^2 n_c
= \alpha\,(b + jc + j^2 a) + \beta\,(c + ja + j^2 b)
= \alpha\,j^2 P + \beta\,j P
= (\alpha j^2 + \beta j)\,P .
\]

\textbf{14.} Avec $j = \frac{-1 + \iu\sqrt3}2$~:
\[
\alpha j = \frac{(3 + \iu\sqrt3)(-1 + \iu\sqrt3)}{12}
= \frac{-3 + 3\iu\sqrt3 - \iu\sqrt3 - 3}{12}
= \frac{-3 + \iu\sqrt3}6 = -\beta ,
\]
donc $\alpha j + \beta = 0$, d'où $\alpha j^2 + \beta j = j(\alpha j
+ \beta) = 0$, et la question 13 donne $n_a + jn_b + j^2n_c = 0$
pour tout triangle $(a, b, c)$. D'après le critère (questions 4
et 5), les centres forment un triangle équilatéral direct ou un
point unique~: c'est le \omterm{pb:b1:complex:1}{théorème de Napoléon}.

\textbf{15.} Le sommet intérieur est $b + \eu^{+\iu\pi/3}(c - b)$,
et le calcul de la question 11, avec $\eu^{+\iu\pi/3} = \frac{1 +
\iu\sqrt3}2$, échange $\alpha$ et $\beta$~: $m_a = \beta b + \alpha
c$, $m_b = \beta c + \alpha a$, $m_c = \beta a + \alpha b$. En
utilisant les identités analogues $b + j^2c + ja = j\,Q$ et $c +
j^2a + jb = j^2 Q$ pour $Q = a + j^2b + jc$~:
\[
m_a + j^2 m_b + j m_c
= \beta(b + j^2 c + ja) + \alpha(c + j^2 a + jb)
= (\beta j + \alpha j^2)\, Q = j(\beta + \alpha j)\,Q = 0 ,
\]
puisque $\alpha j = -\beta$ (question 14). Les centres intérieurs
vérifient donc le critère équilatéral \emph{indirect}~: équilatéral
avec l'orientation opposée (ou réduit à un point).

\textbf{16.} D'après la question 5 appliquée au triplet $(n_a, n_b,
n_c)$ (qui vérifie le critère direct), deux centres coïncident si et
seulement si les trois coïncident~; et les trois coïncident si et
seulement si les \emph{deux} critères sont vérifiés, c'est-à-dire si
et seulement si on a de plus $n_a + j^2 n_b + j n_c = 0$. Le calcul
se mène comme à la question 13, avec les identités $b + j^2c + ja =
jQ$, $c + j^2a + jb = j^2Q$~:
\[
n_a + j^2 n_b + j n_c
= \alpha\,jQ + \beta\,j^2 Q = j(\alpha + \beta j)\,Q .
\]
Directement~: $\beta j = \frac{(3 -
\iu\sqrt3)(-1 + \iu\sqrt3)}{12} = \frac{-3 + 3\iu\sqrt3 +
\iu\sqrt3 + 3}{12} = \frac{\iu\sqrt3}3$, donc $\alpha + \beta j =
\frac{3 + \iu\sqrt3 + 2\iu\sqrt3}6 = \frac{1 + \iu\sqrt3}2 \neq
0$. Le triangle de Napoléon extérieur dégénère donc si et seulement
si $Q = a + j^2b + jc = 0$, c'est-à-dire si et seulement si $(a, b,
c)$ est un triangle équilatéral indirect --- dans ce cas, les
constructions «~vers l'extérieur~» pointent toutes vers l'intérieur
de la région circonscrite au triangle et partagent un même centre.

\textbf{17.} Avec $a = 0$, $b = 1$, $c = \iu$~:
\[
n_a = \alpha + \beta\iu = \frac{(3 + \sqrt3)(1 + \iu)}6,
\qquad
n_b = \alpha\iu = \frac{-\sqrt3 + 3\iu}6,
\qquad
n_c = \beta = \frac{3 - \iu\sqrt3}6 .
\]
Carrés des côtés~: $n_a - n_b = \frac{(3 + 2\sqrt3) +
\iu\sqrt3}6$ donne $\abs{n_a - n_b}^2 = \frac{(3 + 2\sqrt3)^2 +
3}{36} = \frac{24 + 12\sqrt3}{36} = \frac{2 + \sqrt3}3$~; $n_b -
n_c = \frac{(3 + \sqrt3)(-1 + \iu)}6$ donne $\abs{n_b - n_c}^2 =
\frac{2(3 + \sqrt3)^2}{36} = \frac{24 + 12\sqrt3}{36}$~; $n_c - n_a
= \frac{-\sqrt3 - \iu(3 + 2\sqrt3)}6$ redonne la même valeur. Les
trois carrés des côtés valent $\frac{2 + \sqrt3}3$~: le triangle est
équilatéral, comme promis.

\textbf{18.} Développons~:
\[
(a - b)(c - d) + (a - d)(b - c)
= (ac - ad - bc + bd) + (ab - ac - bd + cd)
= ab - ad - bc + cd ,
\]
et $(a - c)(b - d) = ab - ad - bc + cd$~: égalité.

\textbf{19.} On prend les \omterm{def:b1:complex:field}{modules} dans la question 18 et on applique
l'\omterm{prop:b1:complex:rules}{inégalité triangulaire} (\cref{prop:b1:complex:rules}~(4))~:
\[
AC \cdot BD = \abs{(a-b)(c-d) + (a-d)(b-c)}
\leq \abs{a-b}\,\abs{c-d} + \abs{a-d}\,\abs{b-c}
= AB \cdot CD + AD \cdot BC ,
\]
avec égalité si et seulement si les deux termes de la somme sont
portés par une même demi-droite issue de $0$ (le cas d'égalité de
l'\omterm{prop:b1:complex:rules}{inégalité triangulaire}).

\textbf{20.} Factorisation par l'angle moitié
(\cref{met:b1:complex:trig}~(4))~:
$\eu^{\iu\theta} - \eu^{\iu\varphi} = 2\iu\,\sin\frac{\theta -
\varphi}2\;\eu^{\iu(\theta + \varphi)/2}$, et le passage aux \omterm{def:b1:complex:field}{modules}
tue les facteurs de \omterm{def:b1:complex:field}{module} $1$~: $\abs{\eu^{\iu\theta} -
\eu^{\iu\varphi}} = 2\,\abs{\sin\frac{\theta - \varphi}2}$.

\textbf{21.} Avec $p = \eu^{\iu\theta}$, $\theta \in
\intoo{2\pi/3}{4\pi/3}$, la question 20 donne
\[
PA = 2\,\abs{\sin\tfrac\theta2},
\quad
PB = 2\,\abs{\sin\bigl(\tfrac\theta2 - \tfrac\pi3\bigr)},
\quad
PC = 2\,\abs{\sin\bigl(\tfrac\theta2 - \tfrac{2\pi}3\bigr)} .
\]
Sur l'arc, $\frac\theta2 \in \intoo{\pi/3}{2\pi/3}$~: alors
$\sin\frac\theta2 > 0$~; $\frac\theta2 - \frac\pi3 \in
\intoo0{\pi/3}$, donc le deuxième sinus est positif~; $\frac\theta2
- \frac{2\pi}3 \in \intoo{-\pi/3}0$, donc le troisième est négatif
et $PC = 2\sin\bigl(\frac{2\pi}3 - \frac\theta2\bigr)$.
Transformation de somme en produit~:
\[
\sin\Bigl(\frac\theta2 - \frac\pi3\Bigr) +
\sin\Bigl(\frac{2\pi}3 - \frac\theta2\Bigr)
= 2\,\sin\frac\pi6\,\cos\Bigl(\frac\theta2 - \frac\pi2\Bigr)
= \sin\frac\theta2 ,
\]
donc $PB + PC = 2\sin\frac\theta2 = PA$~: c'est le \omterm{pb:b1:complex:1}{théorème de van
Schooten}.

\textbf{22.} En $p = -1$~: $PA = 2$, $PB = PC = 1$ et tous les côtés
du triangle équilatéral ont pour longueur $\sqrt3$, si bien que les
deux membres de Ptolémée valent $2\sqrt3$. Algébriquement, en
utilisant $-1 - j^2 = j$ et $1 + j = -j^2$~:
\[
(a - b)(p - c) = (1 - j)\,(-1 - j^2) = (1 - j)j = j - j^2
= \iu\sqrt3 ,
\]
\[
(a - c)(b - p) = (1 - j^2)(j + 1) = 1 + j - j^2 - j^3
= j - j^2 = \iu\sqrt3 .
\]
Les deux termes sont égaux, donc leur rapport vaut $1 \in
\R_{>0}$~: c'est le cas d'égalité de la question 19, qui traduit
exactement $PA \cdot BC = PB \cdot AC + PC \cdot AB$.

\textbf{23.} Pour $(a, b, c) = (1, j, j^2)$~: $n_a = \alpha j +
\beta j^2 = -\beta + \beta j^2$ (question 14) $= \beta(j^2 - 1)$~;
numériquement, $\beta(j^2 - 1) = \frac{(3 - \iu\sqrt3)}6 \cdot
\bigl(-\frac32 - \iu\frac{\sqrt3}2\bigr) = -1$. De même $n_b =
\alpha j^2 + \beta = -j$ et $n_c = \alpha + \beta j = -j^2$. Le
triangle de Napoléon extérieur de $(1, j, j^2)$ est $(-1, -j,
-j^2)$~: le triangle initial, transformé par la symétrie de centre
son isobarycentre $0$ --- même taille, tourné d'un demi-tour. Un
triangle équilatéral est une forme fixe de la construction de
Napoléon, non une limite qui rétrécit.

\textbf{24.} (i) La multiplication par un nombre de \omterm{def:b1:complex:field}{module} $1$ vue
comme rotation a construit les sommets et les centres (questions 2,
3 et 11) et a converti les distances sur le cercle en sinus
(question 20). (ii) La structure de groupe et l'invariant de forme
ont justifié la normalisation du cercle circonscrit en le cercle
unité et du triangle en $(1, j, j^2)$ à la question 21, via
l'équivariance de la question 10. (iii) La factorisation par
l'angle moitié a alimenté à la fois la question 20 et l'étape de
transformation de somme en produit qui achève van Schooten.

\textbf{25.} Le dictionnaire convertit des énoncés géométriques en
identités polynomiales où chaque hypothèse est une équation~: ce
qu'il apporte, c'est la mécanisation --- le \omterm{pb:b1:complex:1}{théorème de Napoléon}
réduit à $\alpha j + \beta = 0$, une ligne d'arithmétique dans
$\Q(\iu\sqrt3)$~; ce qu'il coûte, c'est la visibilité géométrique
--- le calcul certifie, mais ne montre pas \emph{pourquoi} les
centres se referment en un triangle équilatéral. Une réponse
honnête est que l'identité explique l'\emph{inéluctabilité} du
théorème (elle est vraie identiquement en $a, b, c$, aucune
astuce de configuration n'intervient), tandis que la figure en
explique le \emph{contenu}. Le même dictionnaire revient sous forme
matricielle lorsque les rotations deviennent les matrices
orthogonales $2 \times 2$ du \cref{ch:b1:matrices} et du
\cref{ch:b1:euclid}, où «~la multiplication par $\eu^{\iu\theta}$~»
est le prototype de l'isométrie linéaire.
\end{solution}
